\def\piedfiche{TD n°1 · 1ère S1}
\documentclass[10pt]{article}
\usepackage{amsmath,amssymb}
\usepackage{fontspec}

\setmainfont{Latin Modern Roman}
\newfontfamily\titre{Inter}[UprightFont=* Bold,BoldFont=* ExtraBold]
\usepackage[a4paper,top=1.1cm,bottom=1.35cm,left=1.4cm,right=1.4cm,headheight=0pt,footskip=0.8cm]{geometry}
\usepackage{graphicx}
\usepackage[most]{tcolorbox}
\usepackage{tikz}
\usepackage{enumitem}
\usepackage{fancyhdr}
\usepackage{eso-pic}
\usepackage{multicol}
\usepackage{lastpage}\usepackage{needspace}

% ---------- couleurs Sigma Cap ----------
\definecolor{ink}{HTML}{14284B}
\definecolor{sky}{HTML}{1F8FD0}
\definecolor{lime}{HTML}{6FA82F}
\definecolor{gold}{HTML}{D99A1E}
\definecolor{rouge}{HTML}{C0392B}
\definecolor{violet}{HTML}{6C3FB5}
\definecolor{fond}{HTML}{F2F5F9}

\setlength{\parindent}{0pt}
\setlength{\parskip}{2pt}
\renewcommand{\arraystretch}{1.15}
\newcommand{\R}{\mathbb{R}}
\newcommand{\N}{\mathbb{N}}
\newcommand{\Z}{\mathbb{Z}}
\newcommand{\Q}{\mathbb{Q}}

% ---------- filigrane ----------
\AddToShipoutPictureBG{%
  \AtPageCenter{\makebox(0,0){\rotatebox{50}{\fontsize{78}{78}\selectfont\bfseries\color{ink!6}SIGMA CAP}}}%
  \AtPageLowerLeft{\put(\LenToUnit{0.5\paperwidth},\LenToUnit{0.28\paperheight}){\makebox(0,0){\rotatebox{50}{\fontsize{30}{30}\selectfont\bfseries\color{sky!8}PROFDEMATHS 2.0}}}}%
  \AtPageUpperLeft{\put(0,\LenToUnit{-0.35cm}){\color{ink}\rule{\paperwidth}{0.35cm}}}%
  \AtPageLowerLeft{\put(0,0){\color{lime}\rule{\paperwidth}{0.18cm}}}%
}

% ---------- pied de page ----------
\pagestyle{fancy}\fancyhf{}
\renewcommand{\headrulewidth}{0pt}
\fancyfoot[L]{\scriptsize\color{ink}\textbf{Sigma Cap} -- Vers la réussite · CEM Kennedy}
\fancyfoot[C]{\scriptsize\color{ink!70}\piedfiche}
\fancyfoot[R]{\scriptsize\color{ink}\textbf{\thepage/\pageref{LastPage}}}

% ---------- puces colorées ----------
\setlist[itemize]{leftmargin=1.1em,itemsep=1pt,topsep=2pt,label={\color{sky}\small$\blacktriangleright$}}
\DeclareRobustCommand{\pucenum}[1]{\tikz[baseline=(n.base)]\node[circle,fill=sky,text=white,inner sep=1.3pt,minimum size=1.15em,font=\footnotesize\bfseries](n){#1};}
\setlist[enumerate,1]{leftmargin=1.6em,itemsep=1pt,topsep=1pt,label={\pucenum{\arabic*}}}
\setlist[enumerate,2]{leftmargin=1.5em,itemsep=1pt,topsep=1pt,label={\color{gold}\bfseries\alph*)}}

% ---------- cadres ----------
\newcounter{exo}
\newtcolorbox{exo}[3][]{enhanced,breakable,colback=white,colframe=#2,boxrule=0.9pt,arc=2.5mm,
  left=7pt,right=7pt,top=5pt,bottom=4pt,before skip=3.5pt,after skip=3pt,
  fonttitle=\bfseries\small,coltitle=white,colbacktitle=#2,
  title={\stepcounter{exo}Exercice \theexo\ \ \textcolor{white!85!#2}{\textbullet}\ \ #3},
  attach boxed title to top left={xshift=8pt,yshift=-2.2mm},
  boxed title style={arc=1.5mm,boxrule=0pt},#1}

\newcommand{\partie}[3]{\needspace{7\baselineskip}\vspace{3pt}%
  \begin{tcolorbox}[enhanced,colback=#1,colframe=#1,arc=1.5mm,boxrule=0pt,left=6pt,right=6pt,top=2.5pt,bottom=2.5pt,before skip=4pt,after skip=3pt]
  {\color{white}\bfseries\large #2}\hfill{\color{white}\small\itshape #3}
  \end{tcolorbox}}

\newcommand{\demo}{\tikz[baseline=-0.6ex]\node[fill=violet,text=white,rounded corners=1.5pt,inner sep=1.6pt,font=\scriptsize\bfseries]{DÉMONSTRATION};}
\newcommand{\etoiles}[1]{{\color{gold}\ifcase#1\or$\star$\or$\star\star$\or$\star\star\star$\fi}}

% bulle avec la photo de M. Thiam
\newcommand{\bulle}[1]{%
  \begin{tcolorbox}[enhanced,colback=gold!9,colframe=gold!80!black,boxrule=0.6pt,arc=3mm,
    left=34pt,right=6pt,top=3pt,bottom=3pt,before skip=4pt,after skip=2pt,fontupper=\small,
    overlay={\node[anchor=west] at ([xshift=2pt]frame.west) {\includegraphics[width=26pt]{thiam-rond.png}};}]
  #1\end{tcolorbox}}

\newcommand{\C}{\mathbb{C}}
\newcommand{\entete}[3]{% titre, chapitre, classe
\vspace*{-2pt}
\begin{tcolorbox}[enhanced,colback=fond,colframe=ink,boxrule=1pt,arc=3mm,left=6pt,right=6pt,top=4pt,bottom=4pt]
\begin{minipage}[c]{0.15\textwidth}\centering\includegraphics[height=2.25cm]{logo-full.jpg}\end{minipage}%
\begin{minipage}[c]{0.70\textwidth}\centering
  {\small\color{sky}\bfseries COURS DE RENFORCEMENT · ANNÉE SCOLAIRE 2026 -- 2027}\\[3pt]
  {\LARGE\color{ink}\bfseries #1}\\[2pt]
  {\large\color{ink}\bfseries #2}\\[4pt]
  \tikz\node[fill=lime,text=white,rounded corners=2pt,inner sep=3pt,font=\small\bfseries]{Classe~: #3};\
  \tikz\node[fill=sky,text=white,rounded corners=2pt,inner sep=3pt,font=\small\bfseries]{Prof~: M. Racine Thiam};\
  \tikz\node[fill=gold,text=white,rounded corners=2pt,inner sep=3pt,font=\small\bfseries]{CEM Kennedy};
\end{minipage}%
\begin{minipage}[c]{0.15\textwidth}\centering\includegraphics[height=2.25cm]{profdemaths.jpg}\end{minipage}
\end{tcolorbox}}
\newenvironment{competences}{%
\begin{tcolorbox}[enhanced,colback=sky!6,colframe=sky,boxrule=0.7pt,arc=2mm,left=6pt,right=6pt,bottom=3pt,
  title={\bfseries\small Compétences travaillées dans cette fiche},fonttitle=\small,colbacktitle=sky,coltitle=white,
  attach boxed title to top left={xshift=8pt,yshift=-2mm},boxed title style={arc=1mm,boxrule=0pt},top=6pt]
\small\begin{multicols}{2}\begin{itemize}}{%
\end{itemize}\end{multicols}\vspace{-6pt}
{\footnotesize\color{ink!80}Niveau~: \etoiles1\ application directe \quad \etoiles2\ entraînement \quad \etoiles3\ approfondissement \qquad \demo\ exercice de raisonnement}
\end{tcolorbox}}
\newcommand{\rappel}[1]{\begin{tcolorbox}[enhanced,colback=lime!7,colframe=lime!70!black,boxrule=0.6pt,arc=2mm,left=6pt,right=6pt,top=3pt,bottom=3pt,before skip=3pt,after skip=3pt,fontupper=\small,title={\bfseries\footnotesize Rappel de cours},colbacktitle=lime!70!black,coltitle=white,attach boxed title to top left={xshift=8pt,yshift=-2mm},boxed title style={arc=1mm,boxrule=0pt},top=5pt]#1\end{tcolorbox}}

\begin{document}
\entete{FICHE DE TD N°1}{ÉQUATIONS, INÉQUATIONS, SYSTÈMES ET POLYNÔMES}{1ère S1}

\begin{competences}
  \item Résoudre équations et inéquations du \textbf{second degré}, avec ou sans \textbf{paramètre}.
  \item Se ramener au second degré~: \textbf{changement de variable}, équations \textbf{irrationnelles}, avec \textbf{valeur absolue}.
  \item Utiliser la \textbf{somme} et le \textbf{produit} des racines~; situer un réel par rapport aux racines.
  \item Résoudre des \textbf{systèmes} (déterminant, pivot de Gauss, systèmes non linéaires) et \textbf{mathématiser} un problème.
  \item \textbf{Polynômes}~: racines, factorisation, Horner, division euclidienne, identification.
  \item \textbf{Démontrer}~: formules de Viète, divisibilité, polynômes réciproques, sommes.
\end{competences}

% ================= PARTIE A =================
\partie{ink}{A · Équations et inéquations du second degré}{« Δ est ton meilleur indicateur »}

\begin{exo}{ink}{Échauffement \etoiles1}
Résoudre dans $\R$~:
\[
\text{a) } 2x^2-5x-3=0 \qquad \text{b) } x^2-2\sqrt3\,x+3=0 \qquad \text{c) } x^2-(1+\sqrt2)x+\sqrt2=0 \qquad \text{d) } -3x^2+x-4\geq0
\]
\[
\text{e) } (x^2-1)(-x^2+5x-6)<0 \qquad\qquad \text{f) } \frac{x^2+x-6}{x^2-4x+3}\geq 0 \qquad\qquad \text{g) } \frac{2x-1}{x+2}\leq x-1
\]
\end{exo}

\begin{exo}{ink}{Se ramener au second degré \etoiles2}
Résoudre dans $\R$ en choisissant un \textbf{changement de variable} adapté~:
\[
\text{a) } x^4-5x^2+4=0 \qquad \text{b) } (x^2-3x)^2-2(x^2-3x)-8=0 \qquad \text{c) } x-5\sqrt x+6=0 \qquad \text{d) } \left(\frac{x-1}{x}\right)^2-\frac{x-1}{x}-2=0
\]
\end{exo}

\begin{exo}{ink}{Radicaux et valeurs absolues \etoiles2\ \demo}
Résoudre dans $\R$ en justifiant chaque équivalence~:
\[
\text{a) } \sqrt{x+5}=x-1 \qquad \text{b) } \sqrt{2x-1}<x-2 \qquad \text{c) } \sqrt{x^2-3x}\geq x-1 \qquad \text{d) } |x^2-4|=3x \qquad \text{e) } x^2-|x|-6\leq 0
\]
\end{exo}
\bulle{\textbf{Piège classique~:} $\sqrt A=B \iff \big(B\geq0$ \textbf{et} $A=B^2\big)$. Si tu élèves au carré sans poser $B\geq0$, tu ajoutes de fausses solutions~! Pour $\sqrt A<B$, il faut \textbf{trois} conditions~: $A\geq0$, $B>0$ et $A<B^2$. Et pour $\sqrt A\geq B$, sépare les cas $B<0$ et $B\geq0$.}

\begin{exo}{ink}{Somme et produit des racines \demo\ \etoiles2}
\begin{enumerate}
  \item Soit $ax^2+bx+c=0$ ($a\neq0$) une équation de discriminant $\Delta\geq0$ et de solutions $x'$ et $x''$. À partir des formules des solutions, \textbf{démontrer} que \ $x'+x''=-\dfrac ba$ \ et \ $x'x''=\dfrac ca$.
  \item Sans calculer les solutions $x'$ et $x''$ de \ $3x^2-7x+1=0$, justifier qu'elles existent, puis calculer~:
  \[ x'^2+x''^2\ ;\qquad \frac1{x'}+\frac1{x''}\ ;\qquad x'^3+x''^3\ ;\qquad \frac{x'}{x''}+\frac{x''}{x'}. \]
  \item Former une équation du second degré ayant pour solutions $x'^2$ et $x''^2$.
  \item Trouver deux réels dont la somme vaut $17$ et le produit $72$. Trouver deux réels $u$ et $v$ tels que $u+v=1$ et $uv=1$~: que se passe-t-il~?
\end{enumerate}
\end{exo}

\begin{exo}{ink}{Le grand classique~: équation à paramètre \etoiles3}
Soit $m$ un réel. On considère l'équation \ $(E_m)~:\ (m-1)x^2-2mx+m+2=0$.
\begin{enumerate}
  \item Résoudre $(E_m)$ pour $m=1$.
  \item On suppose $m\neq1$. Montrer que le discriminant réduit vaut $\Delta'=2-m$. Discuter, suivant les valeurs de $m$, l'existence et le \textbf{signe} des solutions de $(E_m)$ (tableau récapitulatif).
  \item Pour quelle valeur de $m$ l'équation admet-elle une solution double~? La calculer.
  \item Déterminer $m$ pour que $0$ soit solution, puis calculer l'autre solution \textbf{sans résoudre} l'équation.
  \item Trouver une \textbf{relation indépendante de $m$} entre la somme $S$ et le produit $P$ des solutions. Retrouver ainsi la solution double.
  \item Déterminer $m$ pour que les solutions vérifient \ $x'^2+x''^2=8$.
  \item Calculer $f_m(1)$ où $f_m(x)=(m-1)x^2-2mx+m+2$. En déduire les valeurs de $m$ pour lesquelles $x'<1<x''$.
  \item Existe-t-il des valeurs de $m$ pour lesquelles $f_m(x)<0$ pour tout réel $x$~?
\end{enumerate}
\end{exo}

\begin{exo}{ink}{Situer un réel par rapport aux racines \demo\ \etoiles3}
Soit $f(x)=ax^2+bx+c$ avec $a\neq0$, et $\alpha$ un réel.
\begin{enumerate}
  \item À l'aide de la forme canonique, démontrer que \ $a\,f(\alpha)=a^2\left(\alpha+\dfrac b{2a}\right)^2-\dfrac{\Delta}{4}$.
  \item En déduire que si $a\,f(\alpha)<0$, alors $\Delta>0$.
  \item On suppose $\Delta>0$. Montrer que $a\,f(\alpha)=a^2(\alpha-x')(\alpha-x'')$, puis que \ $a\,f(\alpha)<0 \iff x'<\alpha<x''$.
  \item \textbf{Application~:} sans calculer $\Delta$, montrer que l'équation $-2x^2+7x+2026=0$ admet deux solutions de signes contraires, et que $3$ est compris entre elles.
\end{enumerate}
\end{exo}
\bulle{\textbf{La règle d'or des paramètres~:} commence \textbf{toujours} par le cas où le coefficient de $x^2$ s'annule (l'équation n'est plus du second degré~!). Ensuite seulement, $\Delta$, puis le signe de $P$ et de $S$. Et pour situer un réel $\alpha$, pense au signe de $a\,f(\alpha)$~: c'est rapide et élégant.}

% ================= PARTIE B =================
\partie{lime!80!black}{B · Systèmes d'équations}{« Mettre en équations, c'est déjà résoudre à moitié »}

\begin{exo}{lime!80!black}{Systèmes linéaires $2\times2$ \etoiles1}
\begin{enumerate}
  \item Résoudre par la méthode des déterminants~: \ $\begin{cases}3x-2y=7\\5x+4y=8\end{cases}$ \quad puis \quad $\begin{cases}\dfrac2x+\dfrac3y=4\\[4pt]\dfrac5x-\dfrac1y=\dfrac32\end{cases}$ \ {\small(poser $X=\frac1x$, $Y=\frac1y$)}.
  \item Discuter suivant le réel $m$ le système \ $\begin{cases}mx+y=2\\x+my=2\end{cases}$ \ (nombre de solutions et solutions).
\end{enumerate}
\end{exo}

\begin{exo}{lime!80!black}{Pivot de Gauss \etoiles2}
Résoudre dans $\R^3$ par la méthode du pivot de Gauss~:
\[
(S_1)~:\ \begin{cases}x+y+z=6\\2x-y+z=3\\x+2y-z=2\end{cases}
\qquad\qquad
(S_2)~:\ \begin{cases}x+2y-z=-3\\2x+y+z=6\\-x+y+2z=3\end{cases}
\]
\end{exo}

\begin{exo}{lime!80!black}{Systèmes non linéaires \etoiles2\ \demo}
\begin{enumerate}
  \item Résoudre~: \ $\begin{cases}x+y=7\\xy=12\end{cases}$ \quad ; \quad $\begin{cases}x^2+y^2=25\\xy=12\end{cases}$ \quad ; \quad $\begin{cases}\dfrac1x+\dfrac1y=\dfrac56\\[4pt]x+y=5\end{cases}$.
  \item \demo\ Démontrer que si $x+y=S$ et $xy=P$, alors $x$ et $y$ sont les solutions de l'équation $X^2-SX+P=0$. À quelle condition sur $S$ et $P$ de tels réels existent-ils~?
\end{enumerate}
\end{exo}

\begin{exo}{lime!80!black}{Mathématiser un problème \etoiles2}
Un groupe d'élèves de Sigma Cap loue un car pour une sortie pédagogique à $60\,000$~F, à partager équitablement. Si $5$ élèves de plus avaient participé, chacun aurait payé $1\,000$~F de moins.
\begin{enumerate}
  \item Choisir l'inconnue (phrase + unité), puis traduire l'énoncé, phrase par phrase, par une équation.
  \item Résoudre et conclure~: combien d'élèves ont participé et combien chacun a-t-il payé~? Vérifier.
\end{enumerate}
\end{exo}
\bulle{\textbf{Mathématiser en 4 temps~:} \textbf{1)} je définis l'inconnue par une phrase (« soit $n$ le nombre d'élèves »)~; \textbf{2)} je traduis chaque phrase~; \textbf{3)} je résous~; \textbf{4)} je vérifie que la solution convient (ici un entier positif~!) et je réponds par une phrase.}

% ================= PARTIE C =================
\partie{sky}{C · Polynômes}{« Une racine trouvée, c'est un facteur gagné »}

\begin{exo}{sky}{Racine évidente, factorisation \etoiles1}
Soit \ $P(x)=2x^3-3x^2-11x+6$.
\begin{enumerate}
  \item Vérifier que $3$ est racine de $P$. Déterminer le polynôme $Q$ tel que $P(x)=(x-3)Q(x)$ \textbf{de deux façons}~: par identification des coefficients, puis par la méthode de Horner.
  \item Factoriser complètement $P(x)$, puis résoudre $P(x)=0$ et $P(x)\leq0$.
\end{enumerate}
\end{exo}

\begin{exo}{sky}{Division euclidienne et divisibilité \etoiles2}
\begin{enumerate}
  \item Effectuer la division euclidienne de \ $A(x)=x^4-3x^3+2x^2+x-1$ \ par \ $B(x)=x^2-x+1$.
  \item Déterminer les réels $a$ et $b$ pour que $f(x)=x^4+1$ soit divisible par $g(x)=x^2+ax+b$. En déduire une factorisation de $x^4+1$ en produit de deux polynômes du second degré.
  \item Déterminer les réels $a$ et $b$ tels que $x^5+ax^4+b$ soit divisible par $(x-1)^2$ \quad {\small(écrire que $1$ est racine, factoriser par $x-1$, puis que $1$ est racine du quotient)}.
\end{enumerate}
\end{exo}

\begin{exo}{sky}{Trois valeurs suffisent \demo\ \etoiles3}
\rappel{Un polynôme de degré inférieur ou égal à $n$ qui admet plus de $n$ racines distinctes est le \textbf{polynôme nul}.}
Soient $a$, $b$, $c$ trois réels distincts. On pose
\[ P(x)=\frac{(x-b)(x-c)}{(a-b)(a-c)}+\frac{(x-a)(x-c)}{(b-a)(b-c)}+\frac{(x-a)(x-b)}{(c-a)(c-b)}. \]
\begin{enumerate}
  \item Quel est le degré maximal de $P$~? Calculer $P(a)$, $P(b)$ et $P(c)$.
  \item En considérant le polynôme $P(x)-1$, démontrer \textbf{sans calcul} que $P(x)=1$ pour tout réel $x$.
  \item \textbf{Prolongement~:} démontrer qu'il existe un unique polynôme $T$ de degré au plus $2$ tel que $T(1)=2$, $T(2)=3$ et $T(3)=6$, et le déterminer.
\end{enumerate}
\end{exo}

\begin{exo}{sky}{Formules de Viète au degré 3 \demo\ \etoiles3}
\begin{enumerate}
  \item On suppose que $x^3+px^2+qx+r=(x-a)(x-b)(x-c)$ pour tout réel $x$. Démontrer par identification que
  \[ a+b+c=-p,\qquad ab+bc+ca=q,\qquad abc=-r. \]
  \item Soit $Q(x)=x^3-2x^2-5x+6$, de racines $a$, $b$, $c$. \textbf{Sans les calculer}, déterminer $a+b+c$, $ab+bc+ca$, $abc$, puis $a^2+b^2+c^2$ et $\dfrac1a+\dfrac1b+\dfrac1c$.
  \item Vérifier ces résultats en trouvant une racine évidente de $Q$ et en factorisant $Q$.
\end{enumerate}
\end{exo}

\begin{exo}{sky}{Polynômes réciproques \etoiles3}
\begin{enumerate}
  \item On considère $(E)~:\ 6x^4-5x^3-38x^2-5x+6=0$. Justifier que $0$ n'est pas solution, puis montrer qu'en divisant par $x^2$ et en posant $X=x+\dfrac1x$, l'équation $(E)$ équivaut à $6X^2-5X-50=0$. Résoudre $(E)$.
  \item Par la même méthode, résoudre \ $2x^4-9x^3+14x^2-9x+2=0$.
  \item \demo\ Soit $P(x)=x^4-13x^3+48x^2-52x+16$. Montrer que pour $x\neq0$, \ $P(x)=x^2\left(X^2-13X+40\right)$ \ avec $X=x+\dfrac4x$. Résoudre $P(x)=0$.
\end{enumerate}
\end{exo}
\bulle{\textbf{Repère les symétries~!} Dans $6x^4-5x^3-38x^2-5x+6$, les coefficients se lisent pareil dans les deux sens~: c'est un polynôme \textbf{réciproque}. Divise par $x^2$, regroupe $x^2+\frac1{x^2}=X^2-2$ et le degré 4 devient un degré 2~!}

% ================= PARTIE D =================
\partie{violet}{D · Démonstrations et problèmes de synthèse}{« Le niveau 1S1, c'est ici »}

\begin{exo}{violet}{Divisibilité par $(x-1)^2$ \demo\ \etoiles3}
Soit $n$ un entier naturel non nul et \ $f_n(x)=nx^{n+1}-(n+1)x^n+1$.
\begin{enumerate}
  \item Démontrer que pour tout réel $x$~: \ $x^n-1=(x-1)\left(x^{n-1}+x^{n-2}+\cdots+x+1\right)$.
  \item Vérifier que $f_n(x)=nx^n(x-1)-(x^n-1)$. En déduire que $f_n(x)=(x-1)\,h_n(x)$ avec $h_n(x)=nx^n-\left(x^{n-1}+\cdots+x+1\right)$.
  \item Calculer $h_n(1)$. En déduire qu'il existe un polynôme $g_n$ tel que $f_n(x)=(x-1)^2g_n(x)$. Déterminer $g_2$.
\end{enumerate}
\end{exo}

\begin{exo}{violet}{Somme des carrés par les polynômes \demo\ \etoiles3}
\begin{enumerate}
  \item Déterminer un polynôme $P$ de degré $3$ tel que $P(0)=0$ et, pour tout réel $x$~: \ $P(x+1)-P(x)=x^2$ \quad {\small(poser $P(x)=ax^3+bx^2+cx$ et identifier)}.
  \item En écrivant cette égalité pour $x=0,1,\dots,n$, démontrer que \ $1^2+2^2+\cdots+n^2=P(n+1)=\dfrac{n(n+1)(2n+1)}{6}$.
  \item Vérifier pour $n=4$, puis calculer $1^2+2^2+\cdots+30^2$.
\end{enumerate}
\end{exo}

\begin{exo}{violet}{Fractions rationnelles \etoiles2}
\begin{enumerate}
  \item Déterminer les réels $a$, $b$, $c$ tels que pour tout $x\neq-2$~: \ $\dfrac{2x^2+3x-1}{x+2}=ax+b+\dfrac{c}{x+2}$.
  \item Déterminer les réels $\alpha$ et $\beta$ tels que pour tout $x\neq\pm3$~: \ $\dfrac{6}{x^2-9}=\dfrac{\alpha}{x-3}+\dfrac{\beta}{x+3}$. En déduire, pour $x\notin\{-3\,;3\}$, une décomposition de $\dfrac{1}{x^4-81}$ \ {\small(écrire $x^4-81=(x^2-9)(x^2+9)$)}.
  \item Résoudre dans $\R$~: \ $\dfrac{1}{x-1}+\dfrac{2}{x+1}\geq\dfrac{4}{x^2-1}$.
\end{enumerate}
\end{exo}
\begin{exo}{violet}{Problème de synthèse~: polynôme à paramètre \demo\ \etoiles3}
Soit $m$ un réel et \ $P_m(x)=x^3+(m-1)x^2+3x-(m+3)$.
\begin{enumerate}
  \item Démontrer que, \textbf{quel que soit} $m$, le réel $1$ est racine de $P_m$.
  \item Déterminer le polynôme $Q_m$ tel que $P_m(x)=(x-1)\,Q_m(x)$.
  \item Calculer le discriminant de $Q_m$ et l'écrire sous forme factorisée. Pour quelle valeur de $m$ le réel $1$ est-il aussi racine de $Q_m$~?
  \item Discuter, suivant les valeurs de $m$, le nombre de solutions de l'équation $P_m(x)=0$. Préciser les solutions pour $m=-2$ et pour $m=6$.
  \item Résoudre $P_6(x)\geq0$.
\end{enumerate}
\end{exo}

\begin{exo}{violet}{Inéquations à paramètre \etoiles3}
\begin{enumerate}
  \item Factoriser $x^2-(m+1)x+m$, puis résoudre, suivant les valeurs du réel $m$, l'inéquation \ $x^2-(m+1)x+m\leq0$.
  \item Déterminer les réels $m$ pour lesquels \ $mx^2+(m-1)x+m-1>0$ \ pour tout réel $x$.
\end{enumerate}
\end{exo}

\bulle{\textbf{Le mot de la fin~:} en 1ère S1, on ne se contente pas de calculer~: on \textbf{justifie}. Écris « on sait que… or… donc… », vérifie tes solutions dans l'équation de départ, et garde un œil sur l'ensemble de définition. Au travail, la réussite t'attend~! \hfill\textbf{M. Thiam}}

\end{document}
